\section*{第 \ref{ch:g6:mass-and-volume} 章 --- 质量与体积}

\begin{solution}{exo:g6:mass-and-volume:1}
体积：物体所占空间的大小 --- 单位为 \unit{cm^3} 和升。
质量：物体所含物质的多少 --- 单位为克和千克。
\end{solution}

\begin{solution}{exo:g6:mass-and-volume:2}
\qty{2000}{cm^3}；\ \qty{350}{cm^3}；\ \qty{1.5}{L}；\
\qty{500}{mL}。
\end{solution}

\begin{solution}{exo:g6:mass-and-volume:3}
读数时把量筒放在桌面上：倾斜的量筒会使液面倾斜。视线：
眼睛偏高或偏低，都会使读到的刻度发生偏移（即斜视错误）。
读弯月面底部：液体沿器壁上升，
只有弯月面的底部才对应真实体积。
\end{solution}

\begin{solution}{exo:g6:mass-and-volume:4}
$265 - 200 = 65$：李子的体积为 \qty{65}{cm^3}。
\end{solution}

\begin{solution}{exo:g6:mass-and-volume:5}
水会贴合石子上的每一处凹凸，因此被挤开的水
的体积恰好等于石子的体积 --- 这正是排水法命题。
直尺需要直边，而水不需要任何形状。
\end{solution}

\begin{solution}{exo:g6:mass-and-volume:6}
$5 \times 2 \times 1 = \qty{10}{cm^3}$。验证：将它放入一个
量筒中 --- 读数应恰好上升 \qty{10}{mL}。
\end{solution}

\begin{solution}{exo:g6:mass-and-volume:7}
$745 - 280 = \qty{465}{g}$ 大米。去皮键：先将空碗放在天平上，
将天平调零，然后再倒入大米。
\end{solution}

\begin{solution}{exo:g6:mass-and-volume:8}
质量相等 --- 各为一千克。羽毛占据的空间大得多，
同样的物质却有大得多的体积。
\end{solution}

\begin{solution}{exo:g6:mass-and-volume:9}
故意让软木塞下沉：用细棒把它压入水中（读取只有软木塞浸没时的
液面上升），或者把它系在一个重物上 --- 先单独测出重物的
液面上升，再测出两者一起的液面上升，
然后减去重物的那部分。
\end{solution}

\begin{solution}{exo:g6:mass-and-volume:10}
倒出共：$3 \times 250 + 200 = \qty{950}{mL}$。声称值与实际
相差 \qty{50}{mL} --- 少了二十分之一升。
\end{solution}

\begin{solution}{exo:g6:mass-and-volume:11}
液面上升 $83.5 - 80.0 = \qty{3.5}{mL}$：项链的总
体积为 \qty{3.5}{cm^3}。水测得的是总体积，包括金属和搭扣在内 ---
因此声称项链含有 \qty{5}{cm^3} 的黄金是不可能的，
因为项链的总体积甚至不足 \qty{5}{cm^3}。
\end{solution}

\begin{solution}{exo:g6:mass-and-volume:12}
直尺给出预测：$4 \times 3 \times 2 = \qty{24}{cm^3}$。将长方体
（用细线系住，不溅水，完全浸没，且不碰到器壁）放入装有 \qty{150}{mL} 水的量筒中，
预测液面将从 \qty{150}{mL} 上升到 \qty{174}{mL}。若水面上升的量
不是 \qty{24}{mL}，超出量筒正常的一两格误差 ---
则排水法命题就将被推翻。
（至今从未被推翻。）
\end{solution}

\begin{solution}{pb:g6:mass-and-volume:1}
\textbf{1.} $20 \times 1000 = \qty{20000}{cm^3}$。
\textbf{2.} $\qty{20}{L} = \qty{2000}{cL}$；$2000 \div 75 = 26$
整瓶（有余数）。
\textbf{3.} $26 \times 75 = \qty{1950}{cL}$，因此 $2000 - 1950 =
\qty{50}{cL}$ 剩余。
\textbf{4.} $50 \div 25 = 2$ 杯，恰好装完。
\textbf{5.} $\qty{75}{cL} = \qty{0.75}{L}$，而一升约重
\qty{1000}{g}：果汁约重 \qty{750}{g}。
\textbf{6.} $750 + 450 = \qty{1200}{g}$ --- 即 \qty{1.2}{kg}。
\textbf{7.} $6 \times 1.2 = \qty{7.2}{kg}$ 瓶子，加上
\qty{2}{kg} 的箱子：共 \qty{9.2}{kg}。
\textbf{8.} $200 \div 9.2$：$21$ 箱重
\qty{193.2}{kg}，$22$ 箱则为 \qty{202.4}{kg} --- 超重。货车
只能装 $21$ 箱。
\textbf{9.} 溢出的是被石子排开的果汁：
石子的体积为 \qty{40}{cm^3}，
由排水法命题得出。
\textbf{10.} $1000 - 40 = \qty{960}{mL}$ 果汁剩余。
\textbf{11.} 量筒的读数应从 \qty{100}{mL} 上升到
\qty{140}{mL} --- 上升 \qty{40}{mL}，与溢出量相符。拳头无法测量，
只有水才说真话。
\textbf{12.} 这 \qty{40}{cm^3} 的空间若装满果汁，
约可容纳 \qty{40}{g} 的果汁 --- 而石子在同样的空间里装入了 \qty{104}{g}
的物质：同样体积下，石子的物质比果汁的物质重得多，
超过两倍。密度一章将给出这一比较
的正式名称和编号。
\end{solution}
